% Compile twice with pdflatex. Bibliography and all macros are included.
\documentclass[11pt,a4paper]{article}
\usepackage[T1]{fontenc}
\usepackage[utf8]{inputenc}
\usepackage{lmodern,amsmath,amssymb,amsthm,mathtools}
\usepackage[margin=28mm,headheight=14pt]{geometry}
\usepackage{microtype,enumitem,booktabs,longtable,array}
\usepackage[dvipsnames]{xcolor}
\usepackage{fancyhdr}
\usepackage{hyperref}
\hypersetup{colorlinks=true,linkcolor=MidnightBlue,citecolor=MidnightBlue,
 urlcolor=MidnightBlue,pdftitle={Singularity of Fuchsian random walks: a writeup of the Lean proof}}
\newcommand{\R}{\mathbb R}
\newcommand{\C}{\mathbb C}
\newcommand{\Z}{\mathbb Z}
\newcommand{\N}{\mathbb N}
\newcommand{\HH}{\mathbb H^2}
\newcommand{\bd}{\partial\HH}
\newcommand{\Ga}{\Gamma}
\newcommand{\ck}{\check\nu}
\newcommand{\supp}{\operatorname{supp}}
\newcommand{\im}{\operatorname{Im}}
\newcommand{\re}{\operatorname{Re}}
\newcommand{\rank}{\operatorname{rank}}
\newcommand{\PSL}{\operatorname{PSL}}
\newcommand{\SL}{\operatorname{SL}}
\newcommand{\id}{\operatorname{id}}
\newcommand{\one}{\mathbf 1}
\newcommand{\dd}{\,d}
\newcommand{\DG}{D_G}
\newcommand{\dG}{d_G}
\newcommand{\Law}{\operatorname{Law}}
\newcommand{\lean}[1]{\par\smallskip\begingroup\raggedright\footnotesize\noindent\textsc{Lean:} \texttt{\detokenize{#1}}.\par\endgroup\smallskip}
\newtheorem{theorem}{Theorem}[section]
\newtheorem{proposition}[theorem]{Proposition}
\newtheorem{lemma}[theorem]{Lemma}
\newtheorem{corollary}[theorem]{Corollary}
\theoremstyle{definition}
\newtheorem{definition}[theorem]{Definition}
\theoremstyle{remark}
\newtheorem{remark}[theorem]{Remark}
\numberwithin{equation}{section}
\setlist{itemsep=3pt,topsep=5pt}
\setlength{\emergencystretch}{3em}
\pagestyle{fancy}\fancyhf{}
\fancyhead[L]{\small Singularity of Fuchsian random walks}
\fancyhead[R]{\small Exposition of the Lean proof}
\fancyfoot[C]{\thepage}
\title{\textbf{Singularity of finitely supported random walks\\on Fuchsian groups}\\[7pt]
\large A mathematical exposition of the Lean formalization}
\author{}
\date{24 September 2026}
\begin{document}
\maketitle
\vspace{-14pt}
\begin{abstract}
We present the argument implemented in the Lean development for singularity
of the hitting measure of a finitely supported, semigroup-generating random
walk on any discrete nonelementary subgroup of $\PSL_2(\R)$. No symmetry or
finite-covolume assumption is made. The cocompact argument combines a
one-sided comparison of invariant currents with an infinite-strip Green
factorization and a Fourier obstruction. For the general case, an ideal end
of a Dirichlet cell supplies a bounded-height endpoint. Boundary minimality
then gives finite separators. Exact entrance identities yield common shadow
estimates, and an explicit Radon--Nikodym rigidity argument forces orbit
quasiconvexity under hypothetical nonsingularity. This reduces to the
cocompact case or a visual-null limit set. The relevant Lean declarations
are indexed throughout and in an appendix.
\end{abstract}
{\small\tableofcontents}
\newpage

\section{Statement, conventions, and the proof's dependencies}
\begin{theorem}\label{thm:main}
Let $\Ga<\PSL_2(\R)$ be a discrete nonelementary subgroup. Let $\mu$ be a
probability measure with finite support $S$, where $S$ generates $\Ga$ as a
semigroup. For independent $\mu$-distributed increments $X_1,X_2,\ldots$,
put $Z_n=X_1\cdots X_n$. Then $Z_no$ converges almost surely to $\bd$, and
its limiting distribution $\nu$ is mutually singular with visual measure:
\[
                         \nu\perp m_o.
\]
This holds for every $o\in\HH$, without symmetry of $\mu$ and without a
restriction on the covolume of $\Ga$.
\end{theorem}
\lean{FuchsianSingularity.fuchsian_hittingMeasure_singular}

Here and below a Lean reference of the form \texttt{Module.declaration}
means the declaration in \texttt{Singularity/Module.lean}; its namespace
is \texttt{Singularity}. The stated scope is finite support with
\emph{semigroup} generation. Group generation alone is not substituted for it.
Nonelementarity means that no point of $\HH\cup\bd$ has a finite orbit.
The upper half-plane has curvature $-1$, and $\bd=\R\cup\{\infty\}$.
At $z=x+iy$ the visual probability is
\begin{equation}\label{eq:poisson}
 dm_z(t)=\frac{y}{\pi((t-x)^2+y^2)}\dd t,\qquad m_z\{\infty\}=0.
\end{equation}
All $m_z$ are equivalent. When $o=i$, its relative Poisson kernel is
\begin{equation}\label{eq:relative-poisson}
 P_z(t)=\frac{dm_z}{dm_i}(t)=\frac{y(1+t^2)}{(t-x)^2+y^2}.
\end{equation}
We use the Busemann sign convention
\[
 \beta_\xi(z,w)=\lim_{u\to\xi}\bigl(d(z,u)-d(w,u)\bigr),\qquad
 P_z(\xi)=e^{\beta_\xi(i,z)}.
\]
In particular $\beta_\infty(z,w)=\log(\im w/\im z)$.

\paragraph{Relation to the formal proof.}
This is a mathematical exposition, not a literal translation of tactic
scripts. It gives the main estimates and their assembly, including the
noncompact argument. We use standard foundational theorems of probability,
measure theory, Fourier analysis, and hyperbolic geometry in their usual
mathematical form. The cocompact Martin--Na\"im boundary package is stated
explicitly in Proposition~\ref{prop:boundary-package}; it is a substantial
standard theorem, not something claimed to follow merely from weak
convergence of hitting laws. The Lean development constructs the needed
package and the other inputs internally; the appendix identifies its
modules. This document does not rederive every foundational theorem or every
one of the development's 2,591 declarations.

The proof is organized without circularity: first prove the cocompact case;
then prove singularity for quasiconvex orbits; next establish the geometric
separator property for noncompact groups; finally derive quasiconvexity
from hypothetical nonsingularity. The final step does not assume a
classification of finitely generated Fuchsian groups, relative hyperbolicity,
or the singularity conclusion of an external rigidity theorem.

\section{The walk, Green functions, and boundary densities}\label{sec:prelim}
Write $e$ for the identity, $p_n(x,y)=\mathbb P_x(Z_n=y)$, and
\begin{equation}\label{eq:green}
 G(x,y)=\sum_{n\ge0}p_n(x,y),\quad G_0=G(e,e),\quad
 F(x,y)=\mathbb P_x(T_y<\infty)=\frac{G(x,y)}{G_0}.
\end{equation}
The equality for $F$ follows by restarting the walk at its first visit to
$y$. Define
\[
 \dG(x,y)=-\log F(x,y),\qquad \DG(g)=\dG(e,g),\qquad D(g)=d(o,go).
\]
The Green distance need not be symmetric. It is left invariant and satisfies
the triangle inequality, since $F(x,z)\ge F(x,y)F(y,z)$.

\begin{proposition}[Basic random-walk facts]\label{prop:basic}
Under the hypotheses of Theorem~\ref{thm:main}:
\begin{enumerate}
\item On counting $\ell^2(\Ga)$,
\[
 (Pf)(x)=\sum_{s\in S}\mu(s)f(xs),\qquad
 (P^*f)(x)=\sum_{s\in S}\mu(s)f(xs^{-1})
\]
satisfy $\|P\|\le1$ and $r(P)<1$. In particular $G=(I-P)^{-1}$ is bounded,
the Green series converges in operator norm, and all $G(x,y)$ are positive.
\item The geometric boundary limit exists almost surely. Its law $\nu$ is
stationary, atomless, supported on the ideal limit set $\Lambda$, and
$g_*\nu\sim\nu$ for every $g\in\Ga$.
\item If $E\subset\bd$ is measurable, then
\begin{equation}\label{eq:martingale}
 (Z_n)_*\nu(E)=\nu(Z_n^{-1}E)\longrightarrow\one_E(Z_\infty)
 \quad\text{almost surely}.
\end{equation}
\item Either $\nu\perp m_o$ or $\nu\ll m_o$.
\end{enumerate}
\end{proposition}
\begin{proof}[Explanation of the inputs and their use]
Jensen's inequality and invariance of counting measure give
$\|Pf\|_2^2\le\sum_s\mu(s)\sum_x|f(xs)|^2=\|f\|_2^2$.
A nonelementary Fuchsian group contains a free subgroup by hyperbolic
ping-pong and is nonamenable. The admissible convolution spectral-radius
criterion gives $r(P)<1$; see Woess~\cite{Woess}, Chapter~12. This is a
spectral-radius assertion and does \emph{not} assert $\|P\|<1$.
The right and left convolution formulations are intertwined by
$f(x)\mapsto f(x^{-1})$. The spectral-radius formula supplies constants
$A<\infty$, $q<1$ with $\|P^n\|\le Aq^n$, proving Green convergence.
Semigroup generation supplies a positive-probability path between any two
states, giving positivity of $G$.

Almost-sure boundary convergence is the standard convergence theorem for
nonelementary walks on $\HH$; in the formal development it is obtained from
the actual path law and the projective action. Stationarity follows by
conditioning on the first increment. Iterating stationarity along positive
words proves equivalence with every group translate. A maximal atom, if
present, would give a nonempty finite invariant set of maximal atoms, which
nonelementarity excludes. Discreteness and escape from finite sets place
the geometric limit in $\Lambda$.

Given the first $n$ increments, the remaining boundary limit has law
$(Z_n)_*\nu$. Thus the left side of \eqref{eq:martingale} is the conditional
expectation of $\one_E(Z_\infty)$; martingale convergence proves the claim.
For the last assertion, take a visual-null Borel set carrying the singular
part of $\nu$ and replace it by the union of all its group translates.
The resulting set $B$ is invariant and visual-null, and
$\nu(B)$ is exactly the mass of the singular part. The event
$Z_\infty\in B$ is invariant under the one-sided shift of the independent
increment sequence: deleting the first increment changes the boundary
limit by a group element. Ergodicity of the independent shift gives
$\nu(B)\in\{0,1\}$. These alternatives give respectively
$\nu\ll m_o$ and $\nu\perp m_o$.
\end{proof}
\lean{FuchsianHittingMeasure; FuchsianBoundaryConcentration; FuchsianHittingErgodicity}

All almost-everywhere identities involving elements of $\Ga$ can be imposed
on one invariant conull set: $\Ga$ is countable and its action is
quasi-invariant. This convention is used throughout.

For any finite quasi-invariant measure $\rho$, put
\begin{equation}\label{eq:cocycle}
 K_\rho(g,\xi)=\frac{d(g_*\rho)}{d\rho}(\xi),\qquad
 \sigma_\rho(g,\xi)=-\log K_\rho(g^{-1},\xi).
\end{equation}
Then
\begin{align}
 \sigma_\rho(gh,\xi)&=\sigma_\rho(g,h\xi)+\sigma_\rho(h,\xi),\label{eq:additive}\\
 \rho(gE)&=\int_E e^{-\sigma_\rho(g,\xi)}\dd\rho(\xi).
 \label{eq:jacobian}
\end{align}
If $\nu=f\rho$ and $H=\log f$, then
\begin{equation}\label{eq:coboundary}
 \sigma_\nu(g,\xi)-\sigma_\rho(g,\xi)=H(\xi)-H(g\xi).
\end{equation}
For $\nu$, $g\mapsto K_\nu(g,\xi)$ is positive $\mu$-harmonic and takes
value one at $e$. Optional stopping at a singleton gives, for any positive
normalized superharmonic $h$,
\begin{equation}\label{eq:superharmonic}
 F(x,y)h(y)\le h(x).
\end{equation}
Applying this in both directions to $K_\nu(\cdot,\xi)$ yields
\begin{equation}\label{eq:cocycle-bound}
 -\dG(g,e)\le\sigma_\nu(g,\xi)\le\DG(g).
\end{equation}
The visual counterpart is $|\sigma_{m_o}(g,\xi)|\le D(g)$.
\lean{StationaryLogCocycle; RadonNikodymCoboundary; VisualLogCocycle}

\section{Green compression and exact separator identities}\label{sec:operators}
These facts apply to any of the groups in the theorem.
\begin{lemma}[Coercivity without laziness]\label{lem:coercive}
For $v\in\ell^2(\Ga)$,
\begin{equation}\label{eq:coercive}
 \re\langle Gv,v\rangle\ge\tfrac12\|v\|_2^2.
\end{equation}
For every $A\subset\Ga$, let $i_A:\ell^2(A)\to\ell^2(\Ga)$ be extension
by zero, and let $r_A=i_A^*$ be restriction. Then
$G_A=r_AGi_A$ is boundedly invertible. Writing $M_A=G_A^{-1}$, we have
$\|M_A\|\le2$.
\end{lemma}
\begin{proof}
Set $w=Gv$, so $v=w-Pw$. Direct expansion gives
\[
 2\re\langle Gv,v\rangle-\|v\|^2
 =\|w\|^2-\|Pw\|^2\ge0.
\]
Extension by zero transfers this inequality to $G_A$. Cauchy--Schwarz
gives $\|G_Au\|\ge\|u\|/2$, so its range is closed. Its adjoint satisfies
the same real-part inequality, hence $\ker G_A^*=\{0\}$; the range is
therefore dense, and thus all of $\ell^2(A)$. The inverse norm bound follows.
No self-adjointness or entrywise positivity of $M_A$ is asserted.
\end{proof}
\lean{Green; GreenAdjoint; Operators}

For $T_A=\inf\{n\ge0:Z_n\in A\}$ define
\[
 F_A(x,a)=\mathbb P_x(T_A<\infty, Z_{T_A}=a),\qquad a\in A.
\]
This counts paths up to their \emph{first} entrance into $A$; by contrast,
$G(x,a)$ counts all visits and places no restriction on intermediate paths.
The first-entrance decomposition gives
\begin{equation}\label{eq:entrance-green}
 G(x,y)=G^{A}(x,y)+\sum_{a\in A}F_A(x,a)G(a,y),
\end{equation}
where $G^{A}$ counts paths killed at entrance into $A$. Since
$0\le F_A(x,a)\le G(x,a)$, positivity implies a bounded entrance operator:
$|F_Au|\le G i_A|u|$. Restricting \eqref{eq:entrance-green} to $y\in A$
gives $Gi_A=F_AG_A$, hence $F_A=Gi_AM_A$.
Consequently, if every supported path from $x$ to $y$ meets $A$,
\begin{equation}\label{eq:separator}
                 G(x,y)=G(x,A)M_AG(A,y).
\end{equation}
Rows and columns of $G$ belong to $\ell^2$ because $G$ is bounded. Thus
the right side is a legitimate Hilbert-space pairing. A two-index matrix
sum should be interpreted through finite-coordinate truncations converging
in that pairing; boundedness of $M_A$ alone does not imply absolute
summability of its entrywise expansion. The equivalent Schur identity is
\[
 G^{A}_{BB}=G_{BB}-G_{BA}G_A^{-1}G_{AB},\qquad B=\Ga\setminus A.
\]
\lean{EntranceOperator; GreenSeparator; FiniteEntranceDecomposition}

\section{The cocompact boundary package}\label{sec:boundary-package}
In this section assume $\Ga\backslash\HH$ is compact. Put
$\check\mu(g)=\mu(g^{-1})$ and let $\ck$ be its hitting law. Then
$\check G(x,y)=G(y,x)$.

\begin{proposition}[Martin kernels and the invariant current]\label{prop:boundary-package}
The geometric circle is the Martin boundary for each of the two walks.
There are positive continuous kernels
\[
 K(g,\eta)=\lim_{y\to\eta}\frac{G(g,y)}{G(e,y)},\qquad
 \check K(g,\xi)=\lim_{x\to\xi}\frac{G(x,g)}{G(x,e)},
\]
which agree almost everywhere with the densities of $g_*\nu$ and
$g_*\ck$. Off the diagonal the limit
\begin{equation}\label{eq:naim}
 \Theta(\xi,\eta)=\lim_{x\to\xi,\,y\to\eta}
       \frac{G(x,y)}{G(x,e)G(e,y)}
\end{equation}
exists, is positive and continuous, and
\begin{equation}\label{eq:current}
           dJ(\xi,\eta)=\Theta(\xi,\eta)\,d\ck(\xi)\,d\nu(\eta)
\end{equation}
is $\Ga$-invariant. The associated flow measure $dJ\,dt$ descends to a
finite, nonzero invariant measure on the compact frame quotient.
\end{proposition}
This is the finite-range cocompact Martin--Ancona--Na\"im package. The
order of the two measures matters: the first is reflected, the second
forward. Background for the boundary results and current convention is
provided by \cite{Gouezel,CantrellTanaka}. The formalization proves the required
Ancona comparison, Martin identification, off-diagonal convergence and
covariance rather than assuming this proposition.

For clarity, the mechanism behind the package is as follows. Cocompactness
gives a uniform orbit net. The spectral bound and exponential hyperbolic
detours make the total weight of paths avoiding a large ball on a geodesic
superexponentially small in the ball radius. First/last-entrance
decompositions, iterated through a summable sequence of relative errors,
give uniform Ancona inequalities
\[
 G(x,y)\asymp G(x,u)G(u,y)
 \quad\text{when }uo\text{ stays a bounded distance from }[xo,yo].
\]
Comparing normalized positive harmonic functions on nested barriers gives
unique Martin limits; integral representations identify them with actual
hitting densities. Off-diagonal strip pairings give the two-variable
limit. Its change-of-basepoint rule is exactly the invariance in
\eqref{eq:current}. Positivity and continuity imply local finiteness of
$J\,dt$ in flow boxes; a compact fundamental cover gives finite mass on
the quotient. These construction statements are the formal results
indexed below, not extra hypotheses imposed on the walk.
\lean{CocompactAxisAncona; GeometricMartinHomeomorph; HittingMartinDensity}
\lean{GeometricNaimKernel; NaimCovariance; GeometricNaimQuotient}

\section{Why one-sided absolute continuity suffices}\label{sec:onesided}
\begin{lemma}\label{lem:onesided}
In the cocompact case, $\nu\ll m_o$ implies that the normalized quotient
measure induced by $J\,dt$ is Liouville probability. In particular,
$\ck\ll m_o$ and $J=cJ_0$, where
\begin{equation}\label{eq:liouville}
                 dJ_0(\xi,\eta)=\frac{d\xi\,d\eta}{(\xi-\eta)^2},\qquad c>0.
\end{equation}
\end{lemma}
\begin{proof}
Let $L$ be Liouville probability on the compact frame quotient and let
$\Omega$ be the normalized probability induced by $J\,dt$. Both are
geodesic-flow invariant. The geodesic flow is ergodic for $L$ by the Hopf
argument (in the formalization, the corresponding Haar assertion follows
from Mautner's argument and the contracting upper/lower unipotent groups).

Fix a continuous observable $f$. For $L$-almost every frame its forward
Ces\`aro averages converge to $\int f\,dL$. This convergence property is
unchanged along a strong stable leaf: the forward trajectories become
asymptotic, and uniform continuity on the compact quotient makes the
difference of their averages tend to zero. It is also unchanged by a
fixed time shift.

In local Hopf coordinates $(\xi^-,\xi^+,t)$, Haar measure is equivalent
to the product of visual measures in the two endpoint coordinates and
Lebesgue measure in the time coordinate. Fubini, followed by stable
saturation and time-shift invariance, says that for visual-almost every
forward endpoint, the convergence property holds for every compatible
backward endpoint and time. The possible diagonal has zero product mass.
Because $\nu\ll m_o$, this full-measure assertion transfers to
$\ck\otimes\nu\otimes dt$, even if $\ck$ is initially singular.
The positive kernel $\Theta$ transfers it further to $\Omega$.

Invariance of $\Omega$ makes the integral of each Ces\`aro average equal
to $\int f\,d\Omega$. Bounded convergence gives
$\int f\,d\Omega=\int f\,dL$. Continuous observables determine finite
measures on the compact quotient, so $\Omega=L$. Lifting this equality to
flow boxes and disintegrating in $t$ gives $J=cJ_0$ on boundary pairs.
Since $\Theta$ is positive, local product rectangles then imply
$\ck\ll m_o$ (and equivalence of both marginals with $m_o$).
\end{proof}
\lean{OneSidedNaimRigidity.geometricNaimQuotientProbability_eq_haar_of_forward_ac}
\lean{StableFrameMeasureTransfer; NaimQuotientComparison}

\begin{lemma}[Uniform densities and Green comparison]\label{lem:green-sharp}
If $\nu$ is nonsingular in the cocompact case, there are positive constants
$a,b,C$ such that
\begin{equation}\label{eq:bounded-density}
 a m_o\le\nu\le b m_o,\qquad a m_o\le\ck\le b m_o,
\end{equation}
and, for all $x,y\in\Ga$,
\begin{equation}\label{eq:sharp}
 C^{-1}e^{-d(xo,yo)}\le G(x,y)\le C e^{-d(xo,yo)}.
\end{equation}
The same comparison holds for $\check G$.
\end{lemma}
\begin{proof}
Proposition~\ref{prop:basic} gives $\nu\ll m_o$, and
Lemma~\ref{lem:onesided} gives $J=cJ_0$. Write the two marginal densities
as $f_-$ and $f_+$. On two separated compact boundary arcs $U,V$, the
identity of currents reads
$\Theta f_-f_+=c\,dJ_0/d(m_o\otimes m_o)$.
Both continuous factors are bounded above and below by positive constants
on $U\times V$. Integrating in one variable bounds the other marginal
density above and below on its arc. The integrals are positive and finite
because the marginals are equivalent probabilities. A finite cover by
such separated rectangles proves \eqref{eq:bounded-density} globally.

First-hit domination gives $F(e,g)g_*\nu\le\nu$. The density bounds imply
$aF(e,g)m_{go}\le b m_o$. The relative Poisson density has essential
supremum $e^{D(g)}$, so $F(e,g)\le(b/a)e^{-D(g)}$.
For the reverse inequality, extend a geodesic through $go$ to a boundary
point $\eta$. Ancona comparison along an orbit sequence on that extension
gives $G(e,g)K(g,\eta)\ge c_0>0$, with $c_0$ uniform. The density
identification and \eqref{eq:bounded-density} give
$K(g,\eta)\le(b/a)e^{D(g)}$. The almost-everywhere bound extends to every
$\eta$ by continuity and full support. Thus $G(e,g)\ge c_1e^{-D(g)}$.
Left invariance gives \eqref{eq:sharp}, and transposition gives the
reflected comparison.
\end{proof}
\lean{NaimCurrentRigidity; GreenVisualUpper; GreenVisualLower}

\section{The periodic strip and its boundary factorization}\label{sec:strip}
Continue to assume that $\Ga\backslash\HH$ is compact and that
$\nu\not\perp m_o$. Choose a hyperbolic element $a\in\Ga$. After
conjugating, its axis is the imaginary axis and $az=e^\tau z$, with
$\tau>0$. We may take $o=i$: changing the basepoint changes visual measure
within its measure class and does not change the abstract walk.
Let $L=\max_{s\in S}d(i,si)$. Choose a closed strip of width greater than
$L$ about the axis and let $A$ consist of the group elements whose orbit
points belong to that strip. Every walk path crossing from its negative
side to its positive side meets $A$.

The strip modulo $\langle a\rangle$ is compact. Properness of the orbit
map, including finite point stabilizers, therefore gives a disjoint
parametrization
\begin{equation}\label{eq:strip}
                 A=\{a^n b_j:n\in\Z,\ 1\le j\le N\}
\end{equation}
with $N<\infty$. Enlarge the width if necessary to ensure $N\ge1$.
The identification in \eqref{eq:strip} identifies $\ell^2(A)$ with
$\ell^2(\Z\times\{1,\ldots,N\})$.

For $\xi<0<\eta$, choose orbit sequences $x_ko\to\xi$ and $y_ko\to\eta$
tracking rays. For sufficiently large $k$, the separator identity gives
\begin{equation}\label{eq:normalized-separator}
 \frac{G(x_k,y_k)}{G(x_k,e)G(e,y_k)}
 =\left(\frac{G(x_k,\cdot)}{G(x_k,e)}\right)_{A}
 M_A
 \left(\frac{G(\cdot,y_k)}{G(e,y_k)}\right)_{A}.
\end{equation}
The vectors on the right converge \emph{strongly in $\ell^2(A)$} to
$(\check K(g,\xi))_{g\in A}$ and $(K(g,\eta))_{g\in A}$, respectively.
Here is the estimate justifying this essential point. The two-sided
Green comparison \eqref{eq:sharp}, followed by the explicit hyperbolic
distance formula, bounds the ratios along a ray to a fixed endpoint on
one side of the axis by
\[
       C_{\xi,j}e^{-|n|\tau}\quad\text{or}\quad
       C_{\eta,j}e^{-|n|\tau},\qquad g=a^n b_j.
\]
The constants are uniform in the point along that ray; bounded tracking
errors change them by a fixed factor. To see the decay at the endpoint
directly, write $b_ji=u_j+iv_j$ and use
\[
 P_{a^n b_ji}(t)=
 \frac{e^{n\tau}v_j(1+t^2)}{(t-e^{n\tau}u_j)^2+e^{2n\tau}v_j^2}
                         =O_t(e^{-|n|\tau}),\qquad t\ne0,\infty.
\]
More explicitly, along a geodesic ray $w_t$ from $i$ to $\xi$,
$d(zo,w_t)-t$ decreases to $\beta_\xi(zo,i)$. Hence
$\exp(d(i,w_t)-d(zo,w_t))\le P_{zo}(\xi)$; a bounded error in tracking
the ray multiplies this bound by a fixed constant. Applying this to
\eqref{eq:sharp} gives the uniform bound just used. Its square is summable over $A$. Coordinatewise
Martin convergence and dominated convergence in $\ell^2$ now apply.
Boundedness of $M_A$, rather than an unjustified interchange of an
entrywise double sum and a boundary limit, yields
\begin{equation}\label{eq:naim-factor}
       \Theta(\xi,\eta)=\check K(A,\xi)^T M_A K(A,\eta).
\end{equation}
The transpose denotes the real bilinear pairing, extended complex
bilinearly when appropriate. All these kernels are real.
\lean{StripNaimKernel; StripNaimContinuity; NaimOrbitPairing}

Write $f_+(x)=d\nu/dx$ and $f_-(x)=d\check\nu/dx$ in the finite boundary
chart. In logarithmic coordinates define
\begin{equation}\label{eq:kplus}
 \begin{split}
 k_g^+(t)&=e^t K(g,e^t)f_+(e^t),\\
 k_g^-(t)&=e^t\check K(g,-e^t)f_-(-e^t).
 \end{split}
\end{equation}
Thus $k_g^+(t)\,dt$ is the restriction of $g_*\nu$ to the positive
half-line, pulled back by $t\mapsto e^t$. The analogous statement for
$k_g^-$ uses $g_*\check\nu$ and $t\mapsto-e^t$. In particular,
\begin{equation}\label{eq:translate-k}
 k_{a^n b_j}^{\pm}(t)=k_{b_j}^{\pm}(t-n\tau),\qquad
 \int_\R k_g^\pm(t)\,dt\le1.
\end{equation}
The density bounds \eqref{eq:bounded-density}, applied after pushforward,
give $k_{b_j}^{\pm}(t)\le C_j e^{-|t|}$. Consequently
\begin{equation}\label{eq:bounded-periodization}
       \sup_t\sum_{n\in\Z}\sum_{j=1}^N k_{a^n b_j}^{\pm}(t)<\infty.
\end{equation}
Define the analysis operators
\[
 H_\pm:L^2(\R)\longrightarrow\ell^2(A),\qquad
        (H_\pm\varphi)_g=\int_\R\varphi(t)k_g^\pm(t)\,dt.
\]
They are bounded: weighted Cauchy--Schwarz gives
\[
 \left|\int\varphi k_g^\pm\right|^2
 \le\left(\int|\varphi|^2 k_g^\pm\right)\left(\int k_g^\pm\right)
 \le\int|\varphi|^2 k_g^\pm,
\]
and summation uses \eqref{eq:bounded-periodization}.
The estimate using $\|k_g^\pm\|_\infty$ separately would not provide this
summable control over the translates.

In the real chart the Liouville current is, up to normalization,
$|\xi-\eta|^{-2}d\xi\,d\eta$. Thus $J=cJ_0$ and
\eqref{eq:naim-factor} give, after absorbing normalization into $c>0$,
\begin{equation}\label{eq:log-kernel}
 c\,\ell(s-t)=k^-(s)^T M_A k^+(t),\qquad
              \ell(u)=\frac1{4\cosh^2(u/2)}.
\end{equation}
All operators now have specified domains. Convolution by $c\ell$ on
$L^2(\R)$, denoted $\mathcal C$, satisfies
\begin{equation}\label{eq:operator-factor}
 \mathcal C=H_-^*M_AH_+:
 L^2(\R)\xrightarrow{H_+}\ell^2(A)
 \xrightarrow{M_A}\ell^2(A)\xrightarrow{H_-^*}L^2(\R).
\end{equation}
Indeed, first pair \eqref{eq:log-kernel} against compactly supported bounded
test functions in both variables. The bound
$\|M_A\|\|k^-(s)\|_2\|k^+(t)\|_2$ and the preceding periodization
estimates justify the vector-valued integrals and their pairings.
Then extend by density and boundedness. This avoids treating $M_A$ as
an entrywise nonnegative matrix, which it need not be.
\lean{KernelFactorization; NormalizedTwoSidedSingularity}

\section{The Fourier obstruction}\label{sec:fourier}
We use $\widehat f(\omega)=\int_\R f(t)e^{-2\pi i\omega t}\,dt$,
with its unitary extension to $L^2$. The following elementary fact is the
finite-generator obstruction for a fixed translation lattice; compare
Bownik~\cite{Bownik}.
\begin{lemma}\label{lem:translates}
For any $h_1,\ldots,h_N\in L^2(\R)$ and $\tau>0$, the closed linear
span of $\{h_j(\cdot-n\tau):n\in\Z,1\le j\le N\}$ is a proper subspace
of $L^2(\R)$.
\end{lemma}
\begin{proof}
Set $q=1/\tau$. For $\omega\in[0,q)$ form the $N\times(N+1)$ matrix
\[
 B(\omega)_{j,l}=\overline{\widehat h_j(\omega+lq)},\qquad 0\le l\le N.
\]
Choose measurable representatives of the Fourier transforms. Rank--nullity
gives $\dim\ker B(\omega)\ge1$ for every $\omega$ where the entries
are finite, hence almost everywhere. We record a measurable choice,
since pointwise existence alone would not suffice.
For $\varepsilon>0$ let
\[
 Q_\varepsilon(\omega)=\varepsilon
       (B(\omega)^*B(\omega)+\varepsilon I)^{-1}.
\]
Its entries are measurable. The spectral theorem for a finite Hermitian
matrix shows that $Q_{1/m}(\omega)$ converges, entry by entry, to the
orthogonal projection $Q(\omega)$ onto $\ker B(\omega)$: on an eigenvalue
$\lambda\ge0$, the multiplier tends to $1$ if $\lambda=0$ and to $0$
otherwise. Thus $Q$ is measurable. Moreover,
\[
 \sum_{l=0}^N\|Q(\omega)e_l\|^2=\operatorname{tr}Q(\omega)\ge1.
\]
At least one column has norm at least $(N+1)^{-1/2}$. Take the first such
column and normalize it, obtaining a measurable unit vector
$v(\omega)\in\ker B(\omega)$.

Define $\widehat\varphi(\omega+lq)=v_l(\omega)$ for $0\le\omega<q$,
$0\le l\le N$, and define it to be zero elsewhere. Then
$\|\widehat\varphi\|_2^2=q>0$. For every $n,j$, Plancherel gives
\[
 \langle\varphi,h_j(\cdot-n\tau)\rangle
 =\int_0^q e^{2\pi in\tau\omega}
       \sum_{l=0}^N v_l(\omega)\overline{\widehat h_j(\omega+lq)}\,d\omega
 =0.
\]
The equality uses $e^{2\pi in\tau lq}=1$. All integrals are legitimate
by Cauchy--Schwarz. This nonzero $\varphi$ proves the assertion.
\end{proof}
\lean{KernelProjection; MeasurableKernel; FrequencyKernel; FourierReduction}

On the other hand, convolution by $\ell$ has trivial kernel. An explicit
calculation compatible with the formal proof is useful. Put
$\gamma(t)=\exp(t-e^t)$ and $\widetilde\gamma(t)=\gamma(-t)$. Substituting
$u=e^s$ gives
\[
 (\gamma*\widetilde\gamma)(t)
 = e^{-t}\int_0^\infty u e^{-(1+e^{-t})u}\,du
 =\frac{e^{-t}}{(1+e^{-t})^2}=\ell(t).
\]
Likewise $\widehat\gamma(\omega)=\Gamma(1-2\pi i\omega)$, so
\begin{equation}\label{eq:liouville-multiplier}
 \widehat\ell(\omega)=|\Gamma(1+2\pi i\omega)|^2
 =\begin{cases}\displaystyle\frac{2\pi^2\omega}{\sinh(2\pi^2\omega)},&\omega\ne0,\\
 1,&\omega=0.
 \end{cases}
\end{equation}
This is everywhere strictly positive. The convolution theorem and
Plancherel imply $\ker\mathcal C=\{0\}$.
\lean{LiouvilleFourier}

Apply Lemma~\ref{lem:translates} to $h_j=k_{b_j}^+$. By
\eqref{eq:translate-k}, it produces a nonzero $\varphi\in\ker H_+$.
Equation~\eqref{eq:operator-factor} would then give
$\mathcal C\varphi=0$, a contradiction. We have proved:
\begin{theorem}\label{thm:cocompact}
Under the hypotheses of Theorem~\ref{thm:main}, if
$\Ga\backslash\HH$ is compact, then $\nu\perp m_o$.
\end{theorem}
\lean{CocompactHittingSingularity; FuchsianCocompactSingularity}

\section{Quasiconvex orbits and a Green-distance criterion}\label{sec:qc}
An orbit $\Ga o$ is \emph{quasiconvex} if some $R<\infty$ satisfies
$[xo,yo]\subset N_R(\Ga o)$ for all $x,y\in\Ga$.
Let $\Lambda\subset\bd$ denote the orbit limit set.
\begin{lemma}\label{lem:qc-dichotomy}
If $\Ga o$ is quasiconvex, either $m_o(\Lambda)=0$ or
$\Ga\backslash\HH$ is compact.
\end{lemma}
\begin{proof}
Suppose $m_o(\Lambda)>0$. Choose a finite-chart Lebesgue density point
$\xi$ of $\Lambda$. Along the vertical ray $w_t=\xi+ie^{-t}$, the
Poisson approximate identity gives $m_{w_t}(\Lambda^c)\to0$.
Quasiconvexity, applied to segments with endpoints tending to $\xi$,
provides $g_t\in\Ga$ at uniformly bounded distance from $w_t$; changing
the initial point of the ray only changes this bound by a constant.
Visual measures at points at distance at most $R'$ have density ratios
between $e^{-R'}$ and $e^{R'}$. Invariance of $\Lambda$ gives
\[
 m_o(\Lambda^c)=m_{g_to}(\Lambda^c)
                 \le e^{R'}m_{w_t}(\Lambda^c)\longrightarrow0.
\]
Since $\Lambda$ is closed and $m_o$ has full support, $\Lambda=\bd$.
Every complete geodesic is then a limit of orbit-endpoint geodesic
segments. Quasiconvexity and properness imply that every point on such
a geodesic lies within a uniform distance of $\Ga o$. Hence the orbit
is coarsely dense in $\HH$. The image of a sufficiently large compact
ball covers the quotient, proving compactness.
\end{proof}
\lean{QuasiconvexConical; QuasiconvexLimitSetDichotomy; QuasiconvexFullBoundary}
Because $\nu(\Lambda)=1$, Lemma~\ref{lem:qc-dichotomy} and
Theorem~\ref{thm:cocompact} prove singularity for every quasiconvex orbit.
\lean{FuchsianQuasiconvexSingularity}

\begin{lemma}\label{lem:green-qc}
If constants $c,C<\infty$ satisfy $\DG(g)\le cD(g)+C$ for every $g$,
then $\Ga o$ is quasiconvex.
\end{lemma}
\begin{proof}
Let $|g|$ be word length for $S\cup S^{-1}$. Since $r(P)<1$, choose
$q<1$ and $A<\infty$ with $\|P^n\|\le Aq^n$. No $S$-path of length
less than $|g|$ reaches $g$, so
\[
 G(e,g)\le\sum_{n\ge |g|}Aq^n=\frac{A}{1-q}q^{|g|}.
\]
Thus $\DG(g)\ge\alpha |g|-\beta$ for some $\alpha>0$.
Combining with the hypothesis gives
$|g|\le c_1D(g)+c_2$. Conversely $D(g)\le L'|g|$, where
$L'=\max_{s\in S\cup S^{-1}}d(o,so)$. By left invariance, the orbit
map is a quasi-isometric embedding of the Cayley graph. Connect adjacent
orbit points by geodesic segments. Images of word-geodesics are uniform
quasi-geodesics; the Morse lemma in $\HH$ puts the geodesic between their
endpoints within a fixed distance of them. Each image edge is itself
within $L'$ of its vertices. This proves orbit quasiconvexity.
\end{proof}
\lean{WordDistance; GreenOrbitQuasiconvex}

\section{Dirichlet ends and bounded-height endpoints}\label{sec:ends}
The remaining argument constructs the finite separators needed to prove
a Green-distance upper bound under hypothetical nonsingularity. It uses
no classification of finitely generated Fuchsian groups.

Call $p\in\bd$ a \emph{bounded-height endpoint} if, for a chart
$B\in\PSL_2(\R)$ with $B\infty=p$, there is $C_B<\infty$ such that
\begin{equation}\label{eq:bounded-height}
                  \im(B^{-1}go)\le C_B\quad(g\in\Ga).
\end{equation}
The condition is independent of the chart: two such charts differ by a
real affine map fixing $\infty$, which multiplies heights by a positive
constant. It is also $\Ga$-invariant, by reindexing $g$.

\begin{lemma}\label{lem:dirichlet-end}
If $\Ga\backslash\HH$ is noncompact, there exists a bounded-height endpoint.
\end{lemma}
\begin{proof}
The closed Dirichlet cell about $o$ is
\[
 \mathcal D=\{w\in\HH:d(w,o)\le d(w,go)\text{ for every }g\in\Ga\}.
\]
Every point has a nearest orbit point, since the orbit is proper.
Translating by that point puts it in $\mathcal D$. If $\mathcal D$ were
bounded, its compact closure would therefore cover the quotient.
Noncompactness supplies $w_n\in\mathcal D$ escaping every ball. A
subsequence converges in the closed disk compactification to $p\in\bd$.
Choose $B\infty=p$ and write $z=B^{-1}o$, $v=B^{-1}go$, $u_n=B^{-1}w_n$.
Then $u_n\to\infty$ and $d(u_n,z)\le d(u_n,v)$.

We claim that $\im v\le\im z$. If $\delta=\im v-\im z>0$, the distance
inequality, using $\cosh d(u,z)=1+|u-z|^2/(2\im u\im z)$, becomes
\[
 \delta|u_n|^2\le2A\re u_n+C,\qquad
 A=(\im v)\re z-(\im z)\re v,\quad
 C=(\im z)|v|^2-(\im v)|z|^2.
\]
This bounds $|u_n|$ by the quadratic inequality
$\delta|u_n|^2\le2|A||u_n|+|C|$, contradicting $u_n\to\infty$.
Thus every $g$ satisfies $\im(B^{-1}go)\le\im(B^{-1}o)$, as required.
\end{proof}
\lean{ProjectiveDirichletCell; DirichletHeightBound; ProjectiveDirichletEnds}

\begin{lemma}\label{lem:dense-endpoints}
For every noncompact quotient, bounded-height endpoints are dense in $\bd$.
\end{lemma}
\begin{proof}
Recall that the closure of every boundary orbit contains $\Lambda$.
For completeness, take a sequence $g_no\to\xi\in\Lambda$ and normalized
matrix representatives converging to a rank-one matrix. Their boundary
actions converge to $\xi$ away from at most one repelling point. Any
nonempty closed invariant set contains a point other than that repelling
point, since nonelementarity excludes finite orbits. It therefore
contains $\xi$. This proves the assertion about orbit closures.

If $\Lambda=\bd$, the orbit of the endpoint supplied by
Lemma~\ref{lem:dirichlet-end} is dense, and bounded height is invariant.
If $\Lambda\ne\bd$, its complement is dense: choose a point outside
$\Lambda$; its orbit stays outside and its closure contains $\Lambda$.
Every $p\notin\Lambda$ is a bounded-height endpoint. Indeed, in a chart
sending $p$ to infinity, an orbit of unbounded height would have a
subsequence converging to infinity, making $p$ an orbit limit point.
This proves density also in this case.
\end{proof}
\lean{BoundedHeightEndpoints; FuchsianDirichletDeficit}

\section{Finite separators and harmonic deficits}\label{sec:deficits}
Take two distinct bounded-height endpoints and send them to $0$ and
$\infty$. The transformed orbit satisfies, for fixed $C,D>0$,
\[
 y\le C,\qquad \frac{y}{x^2+y^2}\le D.
\]
For each $R<\infty$, its intersection with $|x|\le Ry$ lies in a compact
subset of $\HH$, since
\begin{equation}\label{eq:finite-strip}
 \frac1{D(1+R^2)}\le y\le C,\qquad |x|\le RC.
\end{equation}
Consequently only finitely many group elements lie in this axis strip.
Choose $R$ with $\operatorname{arsinh}R>L$, where $L$ bounds jump length.
An edge between opposite sides of the strip cannot avoid it: distance
to the axis is $\operatorname{arsinh}(|x|/y)$, and a crossing geodesic
has to meet the axis. Enlarging the finite strip by the finitely many
one-step predecessors gives a finite set through which every path
exiting the chosen side must pass. In particular, paths whose boundary
limit is outside the closure of that side must visit this finite set.
\lean{BoundedHeightSeparators}

Here is the analytic consequence of a finite separator. For a boundary
region $E$, suppose every path started at $x$ and converging in $E$ visits
a finite set $A$. The strong Markov property at first entrance gives,
as measures restricted to $E$,
\[
                    x_*\nu\big|_E=\sum_{a\in A}F_A(x,a)a_*\nu\big|_E.
\]
Thus, writing $K_\nu(x,\eta)=d(x_*\nu)/d\nu(\eta)$,
\begin{equation}\label{eq:finite-boundary-entrance}
 K_\nu(x,\eta)=\sum_{a\in A}F_A(x,a)K_\nu(a,\eta)
                   \quad\text{for almost every }\eta\in E.
\end{equation}
First-hit domination bounds $K_\nu(a,\eta)\le G_0/G(e,a)$.
Set
\[
 C_A=1+\sum_{a\in A}\frac{G_0}{G(e,a)G(a,e)}.
\]
Then $K_\nu(a,\eta)\le C_AG(a,e)$, and the Green entrance decomposition
implies
\begin{equation}\label{eq:finite-deficit}
 K_\nu(x,\eta)\le C_A\sum_{a\in A}F_A(x,a)G(a,e)
                        \le C_AG(x,e)\quad(\eta\in E\text{ a.e.}).
\end{equation}
For $x=g^{-1}$ this reads
\begin{equation}\label{eq:deficit-bound}
 \DG(g)-\sigma_\nu(g,\eta)\le\log C_A+\log G_0.
\end{equation}
There is no infinite sum or boundary-limit exchange in this argument.
\lean{FiniteEntranceHarmonicDeficit; FuchsianFiniteEntrancePaths}

Define the harmonic shadow and its inverse image by
\begin{equation}\label{eq:harmonic-shadow}
 \mathcal H_R(g)=\{\eta:\DG(g)-\sigma_\nu(g,g^{-1}\eta)\le R\},\qquad
 T_R(g)=g^{-1}\mathcal H_R(g).
\end{equation}
All sets and cocycle formulas may be modified on a single invariant null
set; the group is countable. The cocycle bound \eqref{eq:cocycle-bound}
shows that the deficit in this definition is nonnegative.

\begin{proposition}\label{prop:deficit-compactness}
If bounded-height endpoints are dense, every sequence $(g_n)$ has a
subsequence, still denoted $(g_n)$, and a finite set $Z\subset\bd$ with
the following property: for every open $U\supset Z$ there exists $R$
such that, eventually,
\begin{equation}\label{eq:deficit-compactness}
                         U^c\subset T_R(g_n)\quad\text{modulo a }\nu\text{-null set}.
\end{equation}
One may take $Z$ empty or a singleton.
\end{proposition}
\begin{proof}
If $(g_n)$ has a constant subsequence, its cocycle is bounded below by
\eqref{eq:cocycle-bound}, so take $Z=\varnothing$ and a sufficiently large
$R$. Otherwise pass to a subsequence with $g_n^{-1}o\to\xi\in\bd$ and
take $Z=\{\xi\}$. Given $U\ni\xi$, choose bounded-height endpoints on
either side of $\xi$ so that the closed arc between them containing
$\xi$ is contained in $U$. In the associated chart, $g_n^{-1}o$ is
eventually on that side of the axis. The finite-strip construction
provides a fixed finite separator for every path from these starting
points with boundary limit in $U^c$. Apply \eqref{eq:deficit-bound}.
\end{proof}
\lean{FuchsianDeficitCompactness; FuchsianDirichletDeficit}

\section{Common shadows for harmonic and visual measures}\label{sec:mixed}
Suppose, towards a contradiction, that $\nu\not\perp m_o$. Then
$\nu\ll m_o$ by Proposition~\ref{prop:basic}. There is a measurable,
exactly invariant set $E$ such that
\begin{equation}\label{eq:carrier}
                  \nu\sim\rho:=m_o|_E,\qquad 0<\rho(\bd)\le1.
\end{equation}
Indeed, start with the set where $d\nu/dm_o>0$ and remove the countable
union of its null failures of invariance. Quasi-invariance of both
measures makes this removal harmless. On $E$ the logarithmic cocycle of
$\rho$ equals the visual cocycle. Set $H=\log(d\nu/d\rho)$; it is finite
almost everywhere and satisfies \eqref{eq:coboundary}.
\lean{FuchsianVisualCarrier}

Conjugate so that $o=i$. For $g\in\Ga$ put $t=D(g)$ and choose
$A_g\in\PSL_2(\R)$ with $A_gi=i$ and $A_g(ie^t)=gi$. For $r>0$ define
\[
 \mathcal V_r(g)=A_g\bigl(\bd\setminus[-e^t r,e^t r]\bigr).
\]
For $t=0$, any such choice suffices. Writing $a_tz=e^t z$, the element
$k_g=g^{-1}A_ga_t$ fixes $i$, and
\begin{equation}\label{eq:visual-inverse}
       g^{-1}\mathcal V_r(g)=k_g\bigl(\bd\setminus[-r,r]\bigr).
\end{equation}
Since $k_g$ preserves $m_i$,
\begin{equation}\label{eq:visual-mass}
 m_i\bigl((g^{-1}\mathcal V_r(g))^c\bigr)
                   =\frac2\pi\arctan r\le\frac{2r}\pi.
\end{equation}
The Poisson formula gives, for $\xi\in g^{-1}\mathcal V_r(g)$,
\begin{equation}\label{eq:visual-cocycle}
 D(g)-\log(1+r^{-2})\le\sigma_{m_i}(g,\xi)\le D(g).
\end{equation}
For example, before applying $A_g$ the relevant density at $\eta$ is
$e^t(1+\eta^2)/(e^{2t}+\eta^2)$, which lies between
$e^t/(1+r^{-2})$ and $e^t$ when $|\eta|>e^t r$.
Finally, from every sequence $(g_n)$ we can extract a subsequence with
$k_{g_n}\to k$ in the compact stabilizer of $i$. For any neighborhood
$U$ of $p=k(0)$, choose $r$ small enough; then eventually
\begin{equation}\label{eq:visual-exceptional}
                 (g_n^{-1}\mathcal V_r(g_n))^c\subset U.
\end{equation}
\lean{VisualShadows; VisualShadowExceptionalLimit}

Use a countable increasing family of \emph{mixed} shadows:
\begin{equation}\label{eq:mixed}
 \mathcal S_N(g)=\mathcal H_N(g)\cap\mathcal V_{1/(N+1)}(g),\qquad
               T_N^*(g)=g^{-1}\mathcal S_N(g),\quad N\in\N.
\end{equation}
The following two consequences are the precise shadow geometry needed.
\begin{lemma}\label{lem:mixed-properties}
Assume Proposition~\ref{prop:deficit-compactness}. Then:
\begin{enumerate}
\item There is $N_0$ such that, for each $N\ge N_0$ and
$\lambda\in\{\nu,\rho\}$, constants $c_\lambda>0$ and $C_\lambda<\infty$
independent of $g$ satisfy
\begin{align}
 \lambda(T_N^*(g))&\ge c_\lambda,\label{eq:inverse-positive}\\
 |\sigma_\lambda(g,\xi)-D_\lambda(g)|&\le C_\lambda
                   \quad(\xi\in T_N^*(g)\text{ a.e.}),\label{eq:shadow-cocycle}\\
 e^{-D_\lambda(g)-C_\lambda}\le\lambda(\mathcal S_N(g))
                   &\le e^{-D_\lambda(g)+C_\lambda},\label{eq:shadow-masses}
 \intertext{where}
                  D_\nu(g)&=\DG(g),\qquad D_\rho(g)=D(g).\nonumber
\end{align}
\item From every sequence $(g_n)$ one can extract a subsequence and find
an $N\ge N_0$ and finite $F\subset\Ga$ such that eventually
\begin{equation}\label{eq:finite-cover}
                     \bd=\bigcup_{a\in F}aT_N^*(g_n)
                                      \quad\text{modulo a }\nu\text{-null set}.
\end{equation}
\end{enumerate}
\end{lemma}
\begin{proof}
First, harmonic inverse shadows have uniformly almost full $\nu$-mass
as $N\to\infty$. Otherwise choose $g_N$ with
$\nu(T_N(g_N)^c)\ge\varepsilon$. Apply deficit compactness to a
subsequence. As $\nu$ is atomless, choose an open neighborhood $U$ of
its finite exceptional set with $\nu(U)<\varepsilon$. For some fixed
$R$ and all sufficiently large subsequence indices,
$T_R(g_N)^c\subset U$ modulo null sets. Since $N\ge R$ eventually,
monotonicity contradicts the choice of $g_N$.
Finite-measure absolute continuity transfers this uniform estimate from
$\nu$ to $\rho$. Equation~\eqref{eq:visual-mass} and
$\nu,\rho\ll m_i$ similarly give uniformly almost full mass for the
visual inverse shadows. The intersection therefore has a uniform positive
mass for both measures once $N$ is large.

On the mixed inverse shadow, \eqref{eq:harmonic-shadow} and the global
upper cocycle bound give $\DG(g)-N\le\sigma_\nu(g,\xi)\le\DG(g)$.
Equation~\eqref{eq:visual-cocycle} gives the analogous visual estimate.
Integrating the Jacobian formula \eqref{eq:jacobian} over $T_N^*(g)$,
using its positive mass and the fact that each measure has total mass
at most one, gives \eqref{eq:shadow-masses} after increasing the constants.
This proves (1).

For (2), combine the finite exceptional set from deficit compactness
with the single exceptional point from \eqref{eq:visual-exceptional};
call their union $Z$. There is a finite $F\subset\Ga$ with
$\bigcap_{a\in F}aZ=\varnothing$. Indeed, a point in the intersection
of all translates would have its whole inverse orbit in the finite
set $Z$, contradicting nonelementarity; compactness reduces the empty
intersection to finitely many translates. Choose a sufficiently small
open neighborhood $U$ of $Z$ so that $\bigcap_{a\in F}aU=\varnothing$.
Both exceptional-set statements, and then enlargement of $N$, give
$U^c\subset T_N^*(g_n)$ almost everywhere for all sufficiently large $n$.
Since $\bigcup_{a\in F}aU^c=\bd$ and translates preserve null sets,
\eqref{eq:finite-cover} follows.
\end{proof}
\lean{FuchsianMixedDeficitUniformMass; FuchsianMixedDeficitEstimates; FuchsianDeficitAECover}

\section{From nonsingularity to a uniform Green comparison}\label{sec:rigidity}
We give the measure comparison argument explicitly. This is the step
that turns the preceding geometric information into a global bound;
it is not an application of an external singularity theorem. The use
of common shadows is related to the general Patterson--Sullivan
rigidity framework of Kim--Zimmer~\cite{KimZimmer}.

\begin{proposition}\label{prop:rigidity}
Under \eqref{eq:carrier} and the two conclusions of
Lemma~\ref{lem:mixed-properties}, there is a constant $C<\infty$ such that
\begin{equation}\label{eq:rigid-comparison}
                         |\DG(g)-D(g)|\le C\qquad(g\in\Ga).
\end{equation}
\end{proposition}
\begin{proof}
Choose $M<\infty$ such that the bounded-density window
\[
                         W=\{\xi\in E:|H(\xi)|\le M\}
\]
has positive $\nu$-measure. By \eqref{eq:martingale}, almost every sample
path with limit in $W$ satisfies
\begin{equation}\label{eq:window-concentration}
                  \nu(Z_n^{-1}W^c)\longrightarrow0.
\end{equation}
Since $\rho\ll\nu$ and both are finite, the same is true with $\rho$.
Fix one such path and put $g_n=Z_n$ along it.

Fix any $N\ge N_0$ and write $S_n=\mathcal S_N(g_n)$,
$T_n=g_n^{-1}S_n$. By \eqref{eq:shadow-cocycle},
\eqref{eq:inverse-positive}, and the Jacobian formula,
\[
 \frac{\lambda(S_n\setminus W)}{\lambda(S_n)}
 \le \frac{e^{2C_\lambda}}{c_\lambda}
                     \lambda(g_n^{-1}W^c)\longrightarrow0,
                      \qquad\lambda\in\{\nu,\rho\}.
\]
For all sufficiently large $n$, each shadow therefore has at least
half of its mass in $W$. On $W$, $e^{-M}\rho\le\nu\le e^M\rho$.
It follows that
\[
               \tfrac12e^{-M}\le\frac{\nu(S_n)}{\rho(S_n)}\le2e^M.
\]
Compare this with \eqref{eq:shadow-masses}. There is $L_N<\infty$ with
\begin{equation}\label{eq:path-comparison}
                          |\DG(g_n)-D(g_n)|\le L_N
\end{equation}
for every $n$, after including the finitely many initial values.
The concentration property holds on the chosen path independently of
$N$, so this argument supplies \eqref{eq:path-comparison} for every
integer $N\ge N_0$ on that same path.
\lean{FuchsianShadowPathComparison}

Apply part (2) of Lemma~\ref{lem:mixed-properties} to this sequence.
Pass to its subsequence and fix the resulting $N$ and finite $F$.
For each $a\in F$, quasi-invariance and \eqref{eq:window-concentration}
give
\[
                     \nu(a g_n^{-1}W^c)\longrightarrow0.
\]
Pass to a further subsequence so that these error measures are summable,
simultaneously for all $a\in F$. Borel--Cantelli implies that, for almost
every $\xi$, eventually
\[
                         g_na^{-1}\xi\in W\quad\text{for every }a\in F.
\]
At the same time the finite cover \eqref{eq:finite-cover} supplies
an $a\in F$ with $\eta=a^{-1}\xi\in T_N^*(g_n)$. Consequently both
$\eta\in T_N^*(g_n)$ and $g_n\eta\in W$ hold. The coboundary identity,
the two cocycle estimates, and \eqref{eq:path-comparison} give
\[
 \begin{split}
 |H(\eta)|
 &\le |H(g_n\eta)|+
             |\sigma_\nu(g_n,\eta)-\sigma_\rho(g_n,\eta)|\\
 &\le M+C_\nu+C_\rho+L_N.
 \end{split}
\]
There is also a uniform bound $J$ on
$|\sigma_\nu(a,\zeta)-\sigma_\rho(a,\zeta)|$ for $a\in F$ and almost
every $\zeta$: use the global Green cocycle bounds
\eqref{eq:cocycle-bound} and $|\sigma_\rho(a,\zeta)|\le D(a)$.
Applying \eqref{eq:coboundary} to $\xi=a\eta$ yields
\begin{equation}\label{eq:log-density-bounded}
                 |H(\xi)|\le B:=M+C_\nu+C_\rho+L_N+J
                                    \quad\text{almost everywhere}.
\end{equation}
Thus $e^{-B}\rho\le\nu\le e^B\rho$ globally. This return argument is
also obtainable from vanishing error measures directly, as in the Lean
implementation; the summable subsequence makes the almost-everywhere
conclusion particularly transparent.
\lean{TranslatedConcentrationReturnCover; ReturnCoverDensityRigidity; FuchsianShadowCorrections}

Finally compare the masses of $\mathcal S_N(g)$ for arbitrary $g$, using
\eqref{eq:log-density-bounded} and \eqref{eq:shadow-masses}.
Their ratio is between $e^{-B}$ and $e^B$, hence
\[
                    |\DG(g)-D(g)|\le B+C_\nu+C_\rho.
\]
This proves \eqref{eq:rigid-comparison} for the whole group.
\end{proof}
\lean{BoundedSequenceShadowRigidity; FuchsianShadowGeometryRigidity}

\section{Completion of the proof}\label{sec:completion}
\begin{proof}[Proof of Theorem~\ref{thm:main}]
Boundary convergence and the stated properties of its actual law were
established in Proposition~\ref{prop:basic}. If the quotient is compact,
apply Theorem~\ref{thm:cocompact}.

Suppose the quotient is noncompact and, for a contradiction, that
$\nu\not\perp m_o$. Then $\nu\ll m_o$, and the carrier construction
\eqref{eq:carrier} applies. Lemmas~\ref{lem:dirichlet-end} and
\ref{lem:dense-endpoints} provide a dense set of bounded-height endpoints.
Proposition~\ref{prop:deficit-compactness} supplies deficit compactness;
Lemma~\ref{lem:mixed-properties} supplies the common shadows and their
finite corrected covers. Proposition~\ref{prop:rigidity} now yields
$\DG(g)\le D(g)+C$. Lemma~\ref{lem:green-qc} makes the orbit quasiconvex.
By Lemma~\ref{lem:qc-dichotomy}, either $m_o(\Lambda)=0$ or the quotient
is compact. The first alternative contradicts $\nu\ll m_o$ and
$\nu(\Lambda)=1$; the second contradicts noncompactness.
Thus $\nu\perp m_o$ in all cases.

On $\R$, the density of $m_o$ is strictly positive relative to Lebesgue
measure. The hitting measure has no atom at $\infty$. The same conclusion
therefore holds relative to Lebesgue measure in the finite boundary chart.
\end{proof}
\lean{FuchsianSingularity.fuchsian_hittingMeasure_singular_lebesgue}

The formal final theorem splits the noncompact argument into full and
proper limit sets. The exposition combined them in
Lemma~\ref{lem:dense-endpoints}; the ingredients and resulting separator
estimate are the same. In particular, nonuniform lattices and
infinite-covolume groups both fall within the theorem. Semigroup
generation, finite support, discreteness, and nonelementarity are retained
throughout; symmetry and finite covolume are never added.

\appendix
\section{How to locate the proof in Lean}\label{app:lean}
The entry point is \texttt{Singularity/FuchsianSingularity.lean}.
Its principal declaration is
\begin{quote}\ttfamily\small
Singularity.fuchsian\_hittingMeasure\_singular.
\end{quote}
The group parameter is a subgroup of $\PSL(2,\R)$ with discrete topology.
The law is specified by a finite set $s$, positive weights on $s$, total
mass one, and $\texttt{Submonoid.closure}\,(s)=\top$. The hitting measure
is the law of the constructed geometric boundary limit, not an abstract
measure assumed to have shadow estimates. The companion declaration
\texttt{fuchsian\_randomWalk\_converges\_and\_hittingMeasure\_singular}
includes almost-sure convergence.

The following table lists source modules relative to
\texttt{formalization/Singularity/}; each name has extension
\texttt{.lean}. It is a navigation guide, not a list of additional
hypotheses of the main theorem.
\small
\begin{longtable}{@{}>{\raggedright\arraybackslash}p{0.22\textwidth}>{\raggedright\arraybackslash}p{0.73\textwidth}@{}}
\toprule
Mathematical step & Main source modules\\\midrule
\endfirsthead
\toprule Mathematical step & Main source modules\\\midrule
\endhead
Walk and actual limit & FuchsianHittingMeasure; FuchsianBoundaryConcentration;
FuchsianHittingErgodicity\\[5pt]
Green compression & Green; Operators; GreenSeparator\\[5pt]
Martin--Na\"im package & GeometricNaimKernel; GeometricNaimContinuity;
GeometricNaimCurrent; GeometricNaimQuotient\\[5pt]
One-sided comparison & OneSidedNaimRigidity; NaimCurrentRigidity;
GreenVisualUpper; GreenVisualLower\\[5pt]
Strip factorization & StripNaimKernel; StripNaimContinuity;
KernelFactorization; NormalizedTwoSidedSingularity\\[5pt]
Fourier obstruction & KernelProjection; MeasurableKernel; FrequencyKernel;
FourierReduction; LiouvilleFourier\\[5pt]
Cocompact conclusion & CocompactHittingSingularity; FuchsianCocompactSingularity\\[5pt]
Quasiconvex reduction & GreenOrbitQuasiconvex; QuasiconvexLimitSetDichotomy;
QuasiconvexFullBoundary; FuchsianQuasiconvexSingularity\\[5pt]
Dirichlet end & ProjectiveDirichletCell; DirichletHeightBound;
ProjectiveDirichletEnds; BoundedHeightEndpoints\\[5pt]
Finite separators & BoundedHeightSeparators; FiniteEntranceHarmonicDeficit;
FuchsianFiniteEntrancePaths; FuchsianDirichletDeficit\\[5pt]
Common shadows & FuchsianMixedDeficitShadows; FuchsianMixedDeficitUniformMass;
FuchsianMixedDeficitEstimates; FuchsianDeficitAECover\\[5pt]
Density rigidity & FuchsianShadowPathComparison; TranslatedConcentrationReturnCover;
ReturnCoverDensityRigidity; FuchsianShadowGeometryRigidity\\[5pt]
Final assembly & FuchsianSingularityFromDeficitCompactness;
FuchsianSingularity\\
\bottomrule
\end{longtable}
\normalsize

The recorded project verification reports a successful build of 4471
jobs and an audit of 2591 declarations. The allowed foundational axioms
are \texttt{propext}, \texttt{Classical.choice}, and \texttt{Quot.sound}.
This describes the saved verification report; producing this exposition
does not itself rerun that complete audit. The pinned toolchain is
\texttt{leanprover/lean4:v4.34.0-rc2}, with mathlib commit
\texttt{3649549a1e4b19461e912299ca7127d8831b79fa}.

From the project root, with Lean's \texttt{elan} toolchain manager installed,
the usual reproduction commands are
\begin{quote}\ttfamily
lake exe cache get\\
lake build
\end{quote}
The source distribution includes its toolchain and dependency lock files.
The large \texttt{.lake} directory is a regenerable dependency/build cache.
For the mathematical source of this document, run \texttt{pdflatex} twice
on \texttt{singularity-from-lean.tex}; the bibliography is embedded.

\section{Scope of this exposition and classical inputs}
This document proves the singularity conclusion by the argument organized
in the Lean development, while stating its foundational boundary package
as Proposition~\ref{prop:boundary-package}. It is not a line-by-line
translation of all supporting formal declarations. Standard results used
at that level include the admissible nonamenability spectral criterion,
nonelementary boundary convergence, the cocompact Ancona--Martin theory,
Hopf ergodicity for the compact hyperbolic geodesic flow, the Morse lemma,
Plancherel's theorem, and the martingale convergence theorem. Their
hypotheses and the places where they enter are indicated in the text.
The noncompact Dirichlet-end, finite-entrance, mixed-shadow, and density
rigidity steps are spelled out rather than replaced by an appeal to a
classification or singularity theorem.

A machine-checked source theorem and a human mathematical exposition are
different artifacts: the former certifies its formal statement relative
to its foundations; the latter must also communicate conventions,
dependencies, and mathematical identifications clearly. The source-module
references above are provided to make that comparison possible.

\clearpage
\begin{thebibliography}{9}
\bibitem{Woess}
W.~Woess, \emph{Random Walks on Infinite Graphs and Groups},
Cambridge Tracts in Mathematics \textbf{138}, Cambridge University Press,
2000. \url{https://doi.org/10.1017/CBO9780511470967}.

\bibitem{Gouezel}
S.~Gou\"ezel, \emph{Martin boundary of random walks with unbounded jumps
in hyperbolic groups}, Annals of Probability \textbf{43} (2015),
2374--2404. \url{https://doi.org/10.1214/14-AOP938}.

\bibitem{CantrellTanaka}
S.~Cantrell and R.~Tanaka, \emph{Invariant measures of the topological
flow and measures at infinity on hyperbolic groups}, Journal of Modern
Dynamics \textbf{20} (2024), 215--274.
\url{https://doi.org/10.3934/jmd.2024006}.

\bibitem{Bownik}
M.~Bownik, \emph{The structure of shift-invariant subspaces of
$L^2(\mathbb R^n)$}, Journal of Functional Analysis \textbf{177} (2000),
282--309. \url{https://doi.org/10.1006/jfan.2000.3635}.

\bibitem{KimZimmer}
D.~M.~Kim and A.~Zimmer, \emph{Rigidity for Patterson--Sullivan systems
with applications to random walks and entropy rigidity}, preprint,
arXiv:2505.16556. \url{https://arxiv.org/abs/2505.16556}.
\end{thebibliography}
\end{document}
